\section{The quadratic source and support-aware coding}\label{sec:sparse}

\subsection{The complete parent system}
The common source has 31 positive coordinates, in this exact order:
\[
\begin{gathered}
J,F,\alpha,z,f,h,i,o,j,s,w,\tau,\eta,\zeta,y,\gamma,\\
N,E,P,c,a,X,D,C_2,L,F_2,V,g,I,B_{\rm aux},\Delta.
\end{gathered}
\]
These names are local to this subsection. Put
$q=\beta J+1$, $k=\eta+\zeta$, $H_p=4a+3$ and
$C_p=q-F-\alpha$. The inherited source is the following quadratic system
\cite{package31,package40}.
\begin{longtable}{@{}rl@{}}
\toprule Row & Equation\\\midrule\endhead
1 &$L=q^2-1$\\[2pt]
2 &$N=q(L+1)$\\[2pt]
3 &$X=wq$\\[2pt]
4 &$a=E+sN$\\[2pt]
5 &$E=X(a-E)$\\[2pt]
6 &$c=k(a-E)+\eta$\\[2pt]
7 &$P=E(c-\eta)+k$\\[2pt]
8 &$\tau^2=P(P-k)+1$\\[2pt]
9 &$D=X+ac+\gamma H_p$\\[2pt]
10 &$\Delta=a^2+H_p$\\[2pt]
11 &$C_2=c^2$\\[2pt]
12 &$D^2=\Delta C_2+1$\\[2pt]
13 &$k=jc-V+1+hE$\\[2pt]
14 &$jc-V=(q-F)(N-q)-C_pL+(M_C+qM_F)J$\\[2pt]
15 &$(K_0+w)C_p+q-F-z(q-1)=1$\\[2pt]
16 &$F_2=f^2$\\[2pt]
17 &$V=of-c$\\[2pt]
18 &$g=y^2-V^2$\\[2pt]
19 &$I=iC_2$\\[2pt]
20 &$B_{\rm aux}=I^2$\\[2pt]
21 &$B_{\rm aux}=\Delta(F_2-1)$\\[2pt]
22 &$y^2-B_{\rm aux}g=1$\\
\bottomrule
\end{longtable}
The ten-integer construction uses this table. The nine-variable
construction uses the inherited positive transport shift in row 15,
\[
 (K_0+w)C_p=F+z(q-1).
\]
It corresponds to $z_{\rm old}=z+1$. Canonical compiler witnesses have
$z_{\rm old}\ge2$, so this is an existence equivalence on the positive
domain. Both parent conventions and their shift are checked symbolically;
their monomial supports coincide. They are not treated as literally the
same polynomial system at the same coordinates.

For each parent, replace all coordinates by $x_i+1$, form left-minus-right
residuals, and reorder the rows as
$1,2,3,4,6,7,\ldots,22,5$. Write the result as
$p_1,\ldots,p_{22}\in\Z[x_1,\ldots,x_{31}]$. Each has degree at most two,
and the last has $p_{22}(0)=1$. Thus the all-zero tuple is not a common
zero. The inherited theorem identifies a common natural zero with a
finite ZFC refutation.

\subsection{The weaker separation property}
Pad a monomial of degree less than two by copies of $x_0=1$.
Let $\mathcal S$ be the union of the unordered pairs $(i,j)$ occurring
in the padded supports of the residuals. A weight vector
$w_0=0,w_1,\ldots,w_n\ge0$ has \emph{active-pair isolation} when
\begin{equation}\label{eq:active}
 (i,j)\in\mathcal S,\quad w_i+w_j=w_k+w_l,
 \quad 0\le k\le l\le n
 \quad\Longrightarrow\quad (i,j)=(k,l).
\end{equation}
The comparison is against \emph{all} pairs, not just supported ones.
Requiring uniqueness only among supported pairs would allow an unused
monomial in the square of the packed word to contaminate a tested digit.
Pairs outside $\mathcal S$ may collide with one another.

\begin{theorem}[Support-aware simultaneous extraction]\label{thm:extract}
Let $p_j=\sum_{|\alpha|\le2}a_{j,\alpha}x^\alpha$, $1\le j\le r$,
and suppose \eqref{eq:active} holds. Put $W=\max_iw_i$ and
\[
 E_j=2W+(j-1)(2W+1),\qquad
 c_{j,\alpha}=(2-|\alpha|)!\prod_i\alpha_i!\,a_{j,\alpha}.
\]
For
\[
 \Dcode(Z)=\sum_{j,\alpha}c_{j,\alpha}Z^{E_j-\alpha\cdot w},
 \qquad C_{\mathbf x}(Z)=\left(1+\sum_i x_iZ^{w_i}\right)^2,
\]
all exponents are nonnegative and
\[
 [Z^{E_j}]C_{\mathbf x}(Z)\Dcode(Z)=2p_j(\mathbf x).
\]
If $|c_{j,\alpha}|\le M$, then for $\mathbf x\in\N^n$ every coefficient
of $C_{\mathbf x}\Dcode$ has absolute value at most
$M(1+\sum_i x_i)^2$.
\end{theorem}
\begin{proof}
For a supported $\alpha$, active-pair isolation gives exactly one padded
pair at exponent $\alpha\cdot w$ in $C_{\mathbf x}$. Its coefficient is
$2!x^\alpha/((2-|\alpha|)!\prod_i\alpha_i!)$. The factorial in
$c_{j,\alpha}$ therefore produces $2a_{j,\alpha}x^\alpha$.
A contribution from a different row would equate a nonzero target
difference, of absolute value at least $2W+1$, to a difference of two
pair sums, whose absolute value is at most $2W$. That is impossible.
Within a row, the active monomial is isolated against every possible
pair in $C_{\mathbf x}$, so there are no other contributions.

The reversed supports of the rows lie in disjoint intervals
$[E_j-2W,E_j]$. Each coefficient of $\Dcode$ is consequently bounded by
$M$. All coefficients of $C_{\mathbf x}$ are nonnegative, including
those at collided unused positions, and their sum is
$C_{\mathbf x}(1)=(1+\sum_i x_i)^2$. The coefficient bound follows by
convolution. This last step is why unused collisions cause no hidden
carry problem.
\end{proof}

\subsection{The literal certificate}
In the parent coordinate order above, take
\begin{equation}\label{eq:weightsnew}
\begin{split}
(w_1,\ldots,w_{31})={}&(200,199,193,166,11,57,125,37,93,113,153,16,\\
&158,192,120,64,139,198,144,23,15,43,18,165,179,\\
&112,10,44,117,71,119).
\end{split}
\end{equation}
There are 71 active pairs among the 528 unordered pairs, including
repetitions and the padding coordinate. Exhaustive enumeration verifies
\eqref{eq:active}. The 528 pairs occupy only 285 sums: 130 sum positions
have collisions, and the largest multiplicity is seven. Thus the
condition is strictly weaker than a full Sidon condition in this instance.
The certificate is the literal vector and the exhaustive enumeration;
the heuristic that found it is not a premise of the proof.

Here
\[
 W=200,\quad E_j=400+401(j-1),\quad
 E_{\max}=8821,\quad K_c=E_{\max}+2W=9221.
\]
The previous full-Sidon vector had $W=799$. Both parent conventions have
136 nonzero coefficient records, and the audit verifies all 44 target
identities with the four native coefficient ports left as independent
symbols. The coefficient bound is unchanged, because the source
coefficients and the convolution estimate are unchanged.

This is a reusable sufficient-condition theorem and a specific improved
certificate. We do not assert that support-aware additive coding is new
in all contexts, that $W=200$ is optimal, or that this reduction changes
the undecidability variable-count records.
